% =====================================================================
%  GMN2 §10.5–10.7  $N_f=2$ 的第二种实现 / $N_f=3$ / $N_f=4$：超共形情形
%  原文：wkb-paper.tex 7512–7872 行；PDF（印刷）p.132–140
%  编号公式 (10.18)–(10.30) 共 13 条（其中 (10.20)、(10.29) 原为 multline，
%  (10.23)、(10.26) 原为 align，均改写为 equation+split/aligned，各占一个式号）；
%  插图 7 幅（图 68–74，图号按原文计数硬编码，每处留 % TODO）；脚注 0 条
%  【归属】§10.4「Intermission」（原文 7479–7509 行）不在本块行段内，故未收录
%  （BRIEF 把 7108–7511 行划给 ch14）；若 ch14 也未收，合稿时需补。
%  跨块引用：第 3 节、第 3.2.8 节、第 11 节暂写硬编码并加 % TODO
% =====================================================================

\subsection{$N_f=2$：第二种实现} \label{subsec:Nf2second}

接下来我们考虑 $N_f=2$ 理论的非对称实现：其中 $z=\infty$ 处的不规则奇点与 $N_f=0$ 理论中的相同，而 $z=0$ 处的不规则奇点则被两个正则奇点取代：
\begin{equation}
\lambda^2 = \frac{P_3(z)}{z^2 (z-1)^2} = \left( \frac{m_+^2}{z^2}+\frac{m_-^2}{(z-1)^2} + \frac{\Lambda^2 +u}{2 z} + \frac{\Lambda^2 -u}{2 (z-1)} \right) dz^2
\end{equation}
这里参数 $u,\Lambda$ 未必与上一小节中的相同。绕 $z=\infty$ 的展开显示出一个带单条 WKB 射线的不规则奇点。共有三个转向点，而 $\CB$ 中共有四个奇点。

当质量 $m_\pm$ 非常小时，这里典型的 $T_\WKB$ 有两个退化三角形附着在 $z=\infty$ 处的奇点上。强耦合谱只包含两对 BPS 态，每个奇点一对，如图 68 所示。
% TODO 合稿改 \ref{fig:Nf2evolution-a}
\insfig{5.5in}{Nf2evolution-a}{$N_f=2$ 理论第二种实现的 WKB 三角剖分的演化：取小 $m_\pm$，位于 $u=0$ 附近的强耦合区。为清楚起见，$z=\infty$ 已被映到有限距离处。图中有两对 flip，对应荷为 $\gamma^i_1$、$\gamma^i_2$ 且满足 $\langle \gamma^i_1,\gamma^j_2\rangle=1$ 的两个 BPS 态。注意中间两幅图相差一个 pop，而 pop 并不对应任何 BPS 态。}

与上一节（或第 3 节的分析）相比较可知 $m_\pm = m_1 \pm m_2$。弱耦合谱则表现为一列成对的 BPS 态，它们或与 $z=0$ 处的奇点、或与 $z=1$ 处的奇点相关联（如上一节所述），并产生紧密缠绕的退化三角形。对一般的 $m_\pm$，可能出现退化程度较低的三角剖分，如图 69 所示。
% TODO 上面「第 3 节」合稿改 \ref{sec:physics}
% TODO 合稿改 \ref{fig:Nf2plot-b}
\insfig{2.5in}{Nf2plot-b}{$N_f=2$ 理论（第二种实现）在一般质量下的一个典型 WKB 三角剖分。为清楚起见，$z=\infty$ 已被映到有限距离处。}

$N_f=2$ 理论的两种实现所对应的两个\hk 度量完全相同，这一点绝非显然。事实上，这正是第 3.2.8 节那个出人意料的同构的一个例子。「$N_f=2$ 的两个 Hitchin 系统具有同一个\hk 度量」这一命题的证明如下：这两个系统有相同的谱曲线（尽管它们在 $T^* \IC \IP^1$ 中的嵌入方式不同），也有相同的中心荷函数。现在假设我们能证明（比方说，通过显式考察 $T_\WKB$）：两个系统中的 BPS 谱在某个强耦合点处重合。那么穿墙公式保证两个谱在整个 $\CB$ 上一致；此外，函数 $\CX^\vartheta_\gamma$ 也必然一致，因为它们是同一个黎曼–希尔伯特问题的解，且具有相同的渐近行为。于是这两个 Hitchin 系统的扭量空间与\hk 度量都必定相同。
% TODO 上面「第 3.2.8 节」合稿改 \ref{sec:surprise}

\subsection{$N_f=3$} \label{subsec:Nf3}

带 $N_f=3$ 的 $SU(2)$ 规范理论对应一个 $SU(2)$ Hitchin 系统，它与 $N_f=2$ 理论的第二种实现颇为相似。二次微分为
\begin{equation}
\lambda^2 = \frac{P_4(z)}{z^2 (z-1)^2} = \left(
\frac{m_+^2}{z^2}+\frac{m_-^2}{(z-1)^2} + \frac{2\Lambda m +u}{2 z}
+ \frac{2 \Lambda m -u}{2 (z-1) } + \Lambda^2\right) dz^2
\end{equation}
绕 $z=\infty$ 的展开显示出一个带两条 WKB 射线的不规则奇点，见图 70。
% TODO 合稿改 \ref{fig:Nf3plot}
\insfig{2.5in}{Nf3plot}{带 $N_f=3$ 的 $SU(2)$ 理论的一个典型 WKB 三角剖分。为清楚起见，$z=\infty$ 已被映到有限距离处。}

$z$-平面上有四个转向点，而 $\CB$ 中有五个奇点。三个质量参数 $m$、$m_\pm$ 的地位并不相同。（若要对称地处理它们，就得用 $SU(3)$ Hitchin 系统来实现该理论，这超出了本文的范围。）

当 $m$ 与 $m_\pm$ 都非常小时，$\CB$ 中有四个奇点在 $u=0$ 附近聚拢。在这四个奇点处变为无质量的 BPS 粒子具有相同的规范荷，但味荷不同。见图 71。
% TODO 合稿改 \ref{fig:Nf3evolution-a}
\insfig{5.5in}{Nf3evolution-a}{$N_f=3$ 理论中 $T_\WKB(\vartheta)$ 的演化：取小 $m$ 与 $m_\pm$，位于 $\CB$ 中 $u=0$ 附近的强耦合区。为清楚起见，$z=\infty$ 已被映到有限距离处。随 $\vartheta$ 变化共出现五次 flip，对应荷为 $\gamma^i_1$、$\gamma_2$ 且满足 $\langle \gamma^i_1,\gamma_2\rangle=1$ 的 BPS 态。第二幅与第三幅图相差一个 pop，而 pop 并不对应任何 BPS 态。}

我们知道，当 $m_\pm$ 足够小时，每个正则奇点附近都会有一个转向点；而 BPS 态通常以 $SU(4)$ 味对称的相应 $SU(2)_\pm$ 子群的二重态形式出现。$m$ 则是剩下的那个 $U(1)\in SU(4)$ 的质量参数。这四个粒子实际上构成 $SU(4)$ 的一个 ${\bf 4}$ \cite{Seiberg:1994aj}。的确，在 $SU(2)\times SU(2)\times U(1)$ 之下 ${\bf 4}$ 分解为 $({\bf 2},{\bf 1})_{+1}\oplus({\bf 1},{\bf 2})_{-1}$。从图 71 可以看出，这两个二重态的荷确实相差一个无穷远处的味荷。
% TODO 合稿改 \ref{fig:Nf3evolution-a}

在 $u\sim \frac{\Lambda^2}{4}$ 处还有第二个奇点，那里两个转向点汇合。总而言之，存在一个包含区间 $0<u<\frac{\Lambda^2}{4}$ 的强耦合区，其 BPS 谱很简单：荷为 $\gamma_1^i$ 的粒子的四重态（它们的规范荷全都相同），外加单个荷为 $\gamma_2$、满足 $\langle \gamma_1^i,\gamma_2\rangle=1$ 的粒子；此外对所有 $i,j$ 都有 $\langle \gamma_1^i,\gamma_1^j\rangle=0$。于是穿墙公式就预言了弱耦合区中已知的正确谱——只要把我们先前列出的那些恒等式按一个简单但冗长的次序用一遍，就能看出来。为写出最终结果，令 $\tilde\gamma = 2\gamma_2+\gamma_1^1+\gamma_1^2+\gamma_1^3+\gamma_1^4$。于是：
\begin{equation}
\begin{split}
&\left( \prod_{i=1}^4 \CK_{\gamma_1^i}\right) \CK_{\gamma_2} =
\CK_{\gamma_2} \left( \prod \CK_{\gamma_2+\gamma_1^i}\right) \CK_{\gamma_2+\tilde \gamma}\left( \prod \CK_{\tilde \gamma-\gamma_1^i}\right) \cdots \\
&\quad \cdots \CK_{2\gamma_2+\gamma_1^1+\gamma_1^2+\gamma_1^3+\gamma_1^4}^{-2} \left(\prod_{i<j}\CK_{\gamma_2+\gamma_1^i+\gamma_1^j}\right) \cdots \\
&\quad \cdots \left( \prod \CK_{\tilde \gamma-\gamma_2-\gamma_1^i}\right)\CK_{\tilde \gamma-\gamma_2}\left( \prod_{i=1}^4 \CK_{\gamma_1^i}\right).
\end{split}
\end{equation}

\subsection{$N_f=4$：超共形情形} \label{subsec:SuperConformalCase}

最后，我们已有足够的准备来正经讨论 $N_f=4$ 理论了。本小节只考虑该理论的「平衡」（balanced）实现：$\IC\IP^1$ 上有四个正则奇点和四个转向点。\SW 曲线形如
\begin{equation}
\lambda^2 = \frac{P_4(z)}{D_4(z)^2}(dz)^2= \frac{P^0_4(z)+ u
D_4(z)}{D_4(z)^2}(dz)^2.
\end{equation}
第二个等号中我们引入了一个基点，并定义了一个可归一化的 $u$ 参数。除 $u$ 之外，该理论还有四个质量参数和一个恰好边缘（exactly marginal）的规范耦合。$u$-平面 $\CB$ 中共有 $6$ 个奇点。

若所有质量都取零（或 $u$ 远大于各质量），WKB 流可以变得非常复杂。在这一区域，由 $\lambda$ 决定的 WKB 流可以被下面这个 $\lambda_0$ 的 WKB 流很好地逼近：
\begin{equation}
\lambda_0=\frac{dz}{\sqrt{D_4(z)}} = d\upsilon,
\end{equation}
其中 $\upsilon$ 是椭圆曲线 $y^2=D_4(z)$ 上的单值化坐标（uniformizing coordinate）；于是这些流不过是环面上的直线，其倾角由 $\vartheta$ 决定。对任何有理斜率，我们都得到闭合的 WKB 曲线。特别地，存在一个由 $W$-玻色子构成的无穷谱，它们组成一个 $SL(2,\IZ)$ 对偶轨道。

在 $W$-玻色子出现的同一个 $\vartheta$ 处，还应当出现八个超多重态，它们处于 $Spin(8)$ 味对称群的表示 ${\bf 8}_v$、${\bf 8}_s$ 或 ${\bf 8}_c$ 之中。我们可以这样来理解它们：由于考虑的是非常小的质量，每个转向点都紧挨着四个正则奇点之一。于是，连接两个转向点的 BPS 超多重态会以二重态的形式出现，分别属于两个端点所对应的两个 $SU(2)$ 味子群。若把这四个奇点记作 $a,b,c,d$，把 $Spin(8)$ 中相应的子群记作 $SU(2)_a$、$SU(2)_b$、$SU(2)_c$、$SU(2)_d$，则我们有如下分解
\begin{equation}
\begin{aligned}
{\bf 8}_v &= ({\bf 2}_a, {\bf 2}_b, {\bf 1}_c, {\bf 1}_d) \oplus ({\bf 1}_a, {\bf 1}_b, {\bf 2}_c, {\bf 2}_d), \\
{\bf 8}_s &= ({\bf 2}_a, {\bf 1}_b, {\bf 2}_c, {\bf 1}_d) \oplus ({\bf 1}_a, {\bf 2}_b, {\bf 1}_c, {\bf 2}_d), \\
{\bf 8}_c &= ({\bf 2}_a, {\bf 1}_b, {\bf 1}_c, {\bf 2}_d) \oplus ({\bf 1}_a, {\bf 2}_b, {\bf 2}_c, {\bf 1}_d).
\end{aligned}
\end{equation}
每条闭合 WKB 曲线都把四个奇点分成两对，或是 $(ab,cd)$，或是 $(ac,bd)$，或是 $(ad,bc)$。相应的那对转向点（比方说 $(ab,cd)$）承载着（比方说）${\bf 8}_v$ 中的那 8 个 BPS 超多重态。

一旦打开有限的质量参数，谱就大为简化：每个奇点都变成一个吸引子，使绝大多数闭合 WKB 曲线不再出现。事实上，借助下面的写法引入另一个辅助环面是件有趣的事：
\begin{equation}
\lambda = \frac{dz}{\sqrt{P_4(z)}}\frac{P_4(z)}{D_4(z)}= d\upsilon
\frac{P_4(z)}{D_4(z)},
\end{equation}
其中 $d\upsilon = dz/\sqrt{P_4(z)}$ 现在是环面 $y^2=P_4(z)$ 的单值化参数。WKB 曲线提升为该环面上亚纯微分 $\lambda$ 的流线。$P_4$ 的每个零点在环面上变成一个二重零点，而 $D_4$ 的每个零点则提升为一对极点。若质量参数很小，这对极点会非常靠近那个二重零点。流的局部图像显示出一处局域扰动，它对流而言有一个定义良好的「吸收截面」。我们甚至可以把这个截面算出来：对 $\lambda$ 的局部形式积分，
\begin{equation}
w = \int \lambda \sim \int_{\upsilon=0}^{\upsilon} \frac{m
(\upsilon')^2}{1-(\upsilon')^2}d\upsilon' = - m \upsilon +
\frac{m}{2}\log \frac{1+\upsilon}{1-\upsilon}.
\end{equation}
由于 $- m \upsilon + \frac{m}{2}\log
\frac{1+\upsilon}{1-\upsilon}\sim \frac{m}{3}\upsilon^3 + \cdots$，从 $\upsilon=0$ 发出六条满足 ${\rm Im}(e^{-i \vartheta} w) =0$ 的流。渐近到平行直线的那些流，其间距为 $\pi \mathrm{Re}\left( e^{-i \vartheta}m \right)$。见图 72。
% TODO 合稿改 \ref{fig:perturbed-flow}
\insfig{2.5in}{perturbed-flow}{环面上 $\lambda$ 的一群零点与极点附近的局部流。}

无论质量参数多小，只要一条 WKB 曲线绕环面缠绕足够多次，该奇点最终必定会把它吸住——除非 $\vartheta$ 非常接近质量参数的相位。因此，除非所有质量参数的相位都对齐，看起来在有限的 $u$ 处只有有限多个 $W$-玻色子的 $SL(2,\IZ)$ 像能够存活下来。适当调节质量参数，甚至应当能找到一个没有任何矢量多重态 BPS 态存活的「强耦合区」。那显然会是我们展开分析的好起点。

有鉴于此，让我们考虑下面这个对称的参数选取：
\begin{equation}
\begin{split}
\lambda^2 &= \frac{z^4-u(z^4-1)}{(z^4-1)^2} = \\
&= \frac{1}{16 (z-1)^2}-\frac{1}{16 (z-i)^2}-\frac{1}{16 (z+i)^2}+\frac{1}{16 (z+1)^2}+\frac{1-4 u}{16 (z-1)}\\
&\quad -\frac{i (4 u-1)}{16 (z-i)}+\frac{i (4 u-1)}{16 (z+i)}+\frac{4 u-1}{16 (z+1)}.
\end{split}
\end{equation}
我们把规范耦合调到 $\tau=i$，这给问题带来一个有用的离散 ${\mathbb Z}_4$ 对称性。奇点为 $z_a = 1$、$z_b = i$、$z_c = -1$、$z_d = -i$。质量参数算出来是 $m_a = \frac{1}{4} z_a$。它们的符号是任意的：对这一套参数，我们采用的约定是 $\sum m_a=0$。图 73 给出了一个 WKB 三角剖分的例子。
% TODO 合稿改 \ref{fig:Nf4plot}
\insfig{3.0in}{Nf4plot}{$N_f=4$ 的 $SU(2)$ 理论在特殊 ${\mathbb Z}_4$ 对称点处的一个典型 WKB 三角剖分。}

${\mathbb Z}_4$ 对称性保证 $u$-平面上的六个奇点三三聚拢，在 $u=0$ 与 $u=1$ 处给出两个 $N=4$ 型的奇点。当 $u$ 趋近这两个点时，转向点分别汇合于 $z=0$ 或 $z=\infty$。我们知道，每个 $N=4$ 奇点附近应当有 $6$ 个轻的 BPS 粒子，它们以一切可能的方式连接各转向点（或经由 $z=0$ 附近，或经由 $z=\infty$ 附近）。容易看出，在线段 $0<u<1$ 上，这两组各 $6$ 个粒子就穷尽了 BPS 谱！这一点可以在 $u=1/2$ 处验证；而它在整个线段上都成立，因为中心荷的相位在那里实际上是常数，故不可能发生穿墙。

我们预期会看到一个椭球形的「强耦合」区，其中的谱由这 $12$ 个粒子组成。我们将需要为这些 BPS 粒子的荷取一套方便的标记。首先，我们可以忘掉味荷，模仿 $N=4$ 例子中的标记，投影到一个方便的规范荷格 $\IZ^2$ 上。我们将只遇到荷
\begin{equation}(0,1), (1,1), (1,0), (1,-1). \end{equation}
的确，各转向点排成一个正方形；连接对边的 BPS 弦具有相同的规范荷 $(0,1)$ 或 $(1,0)$，而对角线则具有规范荷 $(1,1)$ 或 $(1,-1)$。

图 74 画出 $T_\WKB$ 的演化，由此可以读出 $u = \half$ 处的 BPS 粒子。
% TODO 合稿改 \ref{fig:Nf4evolution-a}
\insfig{6.2in}{Nf4evolution-a}{$N_f=4$ 的 $SU(2)$ 理论在正文所给参数下 WKB 三角剖分的演化。我们只画出整个演化的一半，即十二次 flip 中的六次。由于问题的 ${\mathbb Z}_4$ 对称性，序列的后一半与前一半相同，只是各图旋转了 $90$ 度。我们又把这一序列分成两行，第一行的最后一幅三角剖分与第二行的第一幅重合。先看到两条对角线边的 flip（对应紫色的 BPS 态），然后是与四个红色 BPS 态中的两个相对应的两条边上的 flip；接着在下一行，先是东北奇点与西南奇点各作一次 pop，随后是与另外两个红色 BPS 态相对应的 flip。演化的另一半包括：由绿色 BPS 态引起的对角线边上的又一次 flip、对应两个蓝色 BPS 态的 flip、西北与东南奇点的 pop，以及由另外两个蓝色 BPS 态带来的又一次 flip。}

共有四个规范荷为 $(0,1)$ 的粒子、四个荷为 $(1,0)$ 的粒子、两个荷为 $(1,1)$ 的粒子和两个荷为 $(1,-1)$ 的粒子，各带不同的味荷组合。让我们暂时忽略味荷信息。这些三角剖分经历一串简单的 flip；其中必然包含每个正则奇点在 $\vartheta=0,\pi$ 或 $\vartheta=\pm\pi/2$ 处所要求的退化三角形。

从图中读出，相应的 KS 变换序列为（这里略去味荷信息，但保留一个 $\pm$ 号，以区分在 $u=0$ 与 $u=1$ 处变轻的粒子）：
\begin{equation}\label{eq:12ParticlesK}
\CK_{-1,1;+}\CK_{1,-1;-}\CK_{0,1;+}^2 \CK_{0,-1;-}^2 \CK_{1,1;+}\CK_{-1,-1;-}\CK_{1,0;+}^2 \CK_{-1,0;-}^2.
\end{equation}

若走到大 $u$ 的区域，各电荷的相位（比如 $(1,0;+)$ 与 $(1,0;-)$）会彼此靠近，因为味荷信息变得无关紧要。这将要求把辛同胚的乘积彻底重排。第一步是把相邻的「边」与「对角线」BPS 态成对地交换（现在 $*$ 表示混合的味荷，即兼有 $+$ 与 $-$ 两种贡献）：
\begin{equation}
\begin{split}
&(\CK_{1,0;-}^2 \CK_{-1,1;+})(\CK_{1,-1;-}\CK_{0,1;+}^2)( \CK_{0,-1;-}^2\CK_{1,1;+})(\CK_{-1,-1;-}\CK_{1,0;+}^2) = \\
&(\CK_{-1,1;+}\CK_{0,1;*}^2\CK_{1,1;*}\CK_{1,0;-}^2 )(\CK_{0,1;+}^2\CK_{1,1;*}\CK_{1,0;*}^2\CK_{1,-1;-}) \\
&(\CK_{1,1;+}\CK_{1,0;*}^2\CK_{1,-1;*} \CK_{0,-1;-}^2)(\CK_{1,0;+}^2\CK_{1,-1;*}\CK_{0,-1;*}^2\CK_{-1,-1;-}).
\end{split}
\end{equation}

我们对左边的四对变换各用了两次五边形恒等式。已经能看出一些有趣的东西：现在有 $8$ 个规范荷为 $(1,0)$ 的粒子、$8$ 个规范荷为 $(0,1)$ 的粒子。我们已经开始看到极大 $u$ 区域中某些超多重态八重态了。另一方面，荷为 $(1,1)$ 的粒子仍然只有 $4$ 个。

朝着大 $u$ 区域走的下一步会产生矢量多重态：在 $\CK_{1,0;-}^2 \CK_{0,1;+}^2$ 这样的成对因子中，我们认出了 $N_f=2$ 穿墙公式的左边；而在 $\CK_{1,-1;-}\CK_{1,1;+}$ 这样的成对因子中，认出了 $N_f=0$ 穿墙公式的左边。后者会产生一个规范荷为 $(2,0)$ 的矢量多重态，前者会产生一个规范荷为 $(2,2)$ 的矢量多重态，并伴随 $4$ 个荷为 $(1,1)$ 的超多重态。它们与上一步已有的另外 $4$ 个超多重态一起，构成一个八重态。

从这以后，穿墙的序列变得越来越复杂。随着 $u$ 增大，各个八重态中每一个成员的相位都朝相应的 $W$-玻色子的相位移动。在这个过程中会遇到由 $N_f=0$ 或 $N_f=2$ 理论的穿墙公式构成的递归结构；把它们在这里反复使用，就产生了 $W$-玻色子在 $SL(2,\IZ)$ S 对偶下的像所构成的无限稠密结构。

从我们这里的讨论反过来读出味结构，是件有意思的事。我们特殊的质量选取把各个 $SU(2)_{a,b,c,d}$ 子群都破缺到它们的嘉当子群，但实际上保留了 $SO(8)$ 中某个并不显然（non-manifest）的部分。在三种可能的 triality 标架中，$SO(8)$ 的嘉当生成元分别取本征值 $(m_a \pm m_c, m_b \pm m_d)$、$(m_a \pm m_b, m_c \pm m_d)$、$(m_a \pm m_d, m_c \pm m_b)$。在第一个标架里，有两个本征值实际上是零：一个隐藏的 $SO(4) \sim SU(2)_L \times SU(2)_R$ 未被破缺。这基本上就是 $N_f=2$ 例子的穿墙公式会在这里重新出现的原因。${\bf 8}_v$ 含这个 $SO(4)$ 的一个矢量与四个单态；${\bf 8}_s$ 与 ${\bf 8}_c$ 则各含两个二重态 ${\bf 2}_L$ 与两个二重态 ${\bf 2}_R$。

我们描述的穿墙，与下面这个看法是相容的：强耦合区中最初那四个荷为 $(1,0)$ 的粒子，应当被视为 ${\bf 8}_s$ 的一部分；更精确地说，可以把 $(1,0;+)$ 取作 $SU(2)_L$ 的一个二重态，把 $(1,0;-)$ 取作 $SU(2)_R$ 的一个二重态。类似地，四个荷为 $(0,1)$ 的粒子在 $8_c$ 中应当是每种二重态各一个：可以把 $(0,1;+)$ 取作 $SU(2)_L$ 的二重态，把 $(0,1;-)$ 取作 $SU(2)_R$ 的二重态。荷为 $(1,1)$ 的粒子（以及同样地 $(1,-1)$ 的粒子）应当全是 $8_v$ 的单态。的确，容易看出第一种穿墙
\begin{equation}
\CK_{1,0;-}^2 \CK_{-1,1;+}=\CK_{-1,1;+}\CK_{0,1;*}^2\CK_{1,1;*}\CK_{1,0;-}^2,
\end{equation}
恰好补上了 ${\bf 8}_s$ 与 ${\bf 8}_c$ 中缺失的二重态以及 ${\bf 8}_v$ 中缺失的单态。而像 $\CK_{1,0;-}^2 \CK_{0,1;+}^2$ 这样把 $SU(2)_L$ 二重态与 $SU(2)_R$ 二重态联系起来的 $N_f=2$ 型穿墙，则补上了 ${\bf 8}_v$ 中缺失的矢量。

描述整个电荷格的一个简单办法如下。在 $u\to 0$ 时变轻的那六个荷（我们标为 $+$）满足与一般 $N=4$ 分析中相同的线性关系，特别地它们张成一个三维格。标为 $-$ 的那六个荷也一样。整个电荷格就是这两个三维格的正交和。

原则上，由强耦合谱与穿墙公式就可以确定参数空间中任何处的谱。下一节（第 11 节）我们将提出一个更好的办法。
% TODO 上面「第 11 节」合稿改 \ref{sec:UniversalStokesMatrix}
